%======================================================================
\section*{Chapter 10: Inferences Based on Two Samples}
\addcontentsline{toc}{section}{Chapter 10: Inferences Based on Two Samples}

\subsection*{10.1 The Two-Sample $z$ CI and Test}
\addcontentsline{toc}{subsection}{10.1 The Two-Sample $z$ CI and Test}

\textbf{Assumptions.}
\begin{enumerate}
  \item $X_1,\dots,X_m$ is a random sample from a population with mean $\mu_1$ and standard deviation $\sigma_1$.
  \item $Y_1,\dots,Y_n$ is a random sample from a population with mean $\mu_2$ and standard deviation $\sigma_2$.
  \item The $X$ and $Y$ samples are independent of each other.
\end{enumerate}

Consider the parameter $\mu_1-\mu_2$. The natural estimator of $\mu_1-\mu_2$ is $\bar{X}-\bar{Y}$;
the test statistic is the standardized form of $\bar{X}-\bar{Y}$. We have
\[
E(\bar{X}-\bar{Y}) = E(\bar{X})-E(\bar{Y}) = \mu_1-\mu_2,
\]
\[
V(\bar{X}-\bar{Y}) = V(\bar{X})+V(\bar{Y}) = \frac{V(X_1+\dots+X_m)}{m^2} + \frac{V(Y_1+\dots+Y_n)}{n^2} = \frac{\sigma_1^2}{m}+\frac{\sigma_2^2}{n},
\]
\[
\sigma_{\bar{X}-\bar{Y}} = \sqrt{\frac{\sigma_1^2}{m}+\frac{\sigma_2^2}{n}}.
\]

\textbf{Proposition.} The expectation of $\bar{X}-\bar{Y}$ is $\mu_1-\mu_2$, so $\bar{X}-\bar{Y}$ is an unbiased estimator of $\mu_1-\mu_2$. Its standard deviation is
$\sqrt{\sigma_1^2/m+\sigma_2^2/n}$.

With $\theta=\mu_1-\mu_2$, $\hat\theta=\bar{X}-\bar{Y}$ and $\sigma_{\hat\theta}=\sigma_{\bar{X}-\bar{Y}}$, the test statistic value is
\[
\frac{\hat\theta - \text{null value}}{\sigma_{\hat\theta}} .
\]

\textbf{The two-sample $z$ CI for $\mu_1-\mu_2$ with known $\sigma$'s.}
Assumptions: (1) the population distributions are normal; (2) $\sigma_1$ and $\sigma_2$ are known.
The pivotal quantity is
\[
Z = \frac{\bar{X}-\bar{Y}-(\mu_1-\mu_2)}{\sqrt{\dfrac{\sigma_1^2}{m}+\dfrac{\sigma_2^2}{n}}} \sim N(0,1).
\]
Then we have the probability statement
\[
P\!\left(-z_{\alpha/2} < \frac{\bar{X}-\bar{Y}-(\mu_1-\mu_2)}{\sigma_{\bar{X}-\bar{Y}}} < z_{\alpha/2}\right) = 1-\alpha
\]
\[
\Rightarrow\; P\!\left(\bar{X}-\bar{Y}-z_{\alpha/2}\,\sigma_{\bar{X}-\bar{Y}} < \mu_1-\mu_2 < \bar{X}-\bar{Y}+z_{\alpha/2}\,\sigma_{\bar{X}-\bar{Y}}\right) = 1-\alpha .
\]
This implies the $100(1-\alpha)\%$ CI for $\mu_1-\mu_2$ with known $\sigma$'s: for two independent samples from normal
population distributions, the $100(1-\alpha)\%$ CI has endpoints
\[
\bar{x}-\bar{y} \pm z_{\alpha/2}\,\sigma_{\bar{X}-\bar{Y}},
\]
with lower limit $\bar{x}-\bar{y}-z_{\alpha/2}\sigma_{\bar{X}-\bar{Y}}$ and upper limit $\bar{x}-\bar{y}+z_{\alpha/2}\sigma_{\bar{X}-\bar{Y}}$.
The corresponding one-sided confidence bounds are
$\bar{x}-\bar{y}-z_\alpha\sigma_{\bar{X}-\bar{Y}}$ (lower bound) and $\bar{x}-\bar{y}+z_\alpha\sigma_{\bar{X}-\bar{Y}}$ (upper bound).

\textbf{Sample size determination.} For known $\sigma_1$ and $\sigma_2$ we can assume $m=n$. The CI width is
\[
W = 2z_{\alpha/2}\sqrt{\frac{\sigma_1^2}{m}+\frac{\sigma_2^2}{n}}
\;\overset{m=n}{=}\; 2z_{\alpha/2}\sqrt{\frac{\sigma_1^2+\sigma_2^2}{n}}
\;\Rightarrow\; n = \frac{z_{\alpha/2}^2(\sigma_1^2+\sigma_2^2)}{(W/2)^2}.
\]

\textbf{Test procedures for $\mu_1-\mu_2$ with known $\sigma$'s (two-sample $z$ test).}
Null hypothesis $H_0:\mu_1-\mu_2=\Delta_0$. Test statistic value:
\[
z = \frac{\bar{x}-\bar{y}-\Delta_0}{\sqrt{\dfrac{\sigma_1^2}{m}+\dfrac{\sigma_2^2}{n}}}.
\]

\begin{center}
\begin{tabular}{lll}
\toprule
Alternative hypothesis & Rejection region & $P$-value\\
\midrule
$H_a:\mu_1-\mu_2>\Delta_0$   & $z\ge z_\alpha$ (upper-tailed) & $1-\Phi(z)$\\
$H_a:\mu_1-\mu_2<\Delta_0$   & $z\le -z_\alpha$ (lower-tailed) & $\Phi(z)$\\
$H_a:\mu_1-\mu_2\ne\Delta_0$ & either $z\le -z_{\alpha/2}$ or $z\ge z_{\alpha/2}$ (two-tailed) & $2[1-\Phi(|z|)]$\\
\bottomrule
\end{tabular}
\end{center}
These concepts may be applied to nonnormal distributions when $m,n>40$ (by the CLT).

\textbf{$\beta$, power, and sample size determination.}
Consider $H_a:\mu_1-\mu_2>\Delta_0$ with rejection region $z\ge z_\alpha$:
\[
z = \frac{\bar{x}-\bar{y}-\Delta_0}{\sqrt{\sigma_1^2/m+\sigma_2^2/n}} \ge z_\alpha
\;\Leftrightarrow\; \bar{x}-\bar{y}\ge \Delta_0 + z_\alpha\sqrt{\frac{\sigma_1^2}{m}+\frac{\sigma_2^2}{n}} .
\]
So $H_0$ is not rejected when $\bar{x}-\bar{y}<\Delta_0+z_\alpha\sigma_{\bar{X}-\bar{Y}}$.
Let $\Delta'$ be a particular value of $\mu_1-\mu_2$ that exceeds $\Delta_0$. Then
\begin{align*}
\beta(\Delta') &= P(\text{type II error}) = P(H_0 \text{ is not rejected when it is false})\\
&= P\!\left(\bar{X}-\bar{Y}<\Delta_0+z_\alpha\sigma_{\bar{X}-\bar{Y}} \text{ when } \mu_1-\mu_2=\Delta'\right)\\
&= P\!\left(\frac{\bar{X}-\bar{Y}-\Delta'}{\sigma_{\bar{X}-\bar{Y}}} < \frac{\Delta_0-\Delta'}{\sigma_{\bar{X}-\bar{Y}}}+z_\alpha \text{ when } \mu_1-\mu_2=\Delta'\right)
= \Phi\!\left(z_\alpha - \frac{\Delta'-\Delta_0}{\sigma_{\bar{X}-\bar{Y}}}\right).
\end{align*}
Similarly, with $\sigma_{\bar{X}-\bar{Y}}=\sqrt{\sigma_1^2/m+\sigma_2^2/n}$:

\begin{center}
\begin{tabular}{ll}
\toprule
Alternative hypothesis & $\beta(\Delta')$\\
\midrule
$H_a:\mu_1-\mu_2>\Delta_0$   & $\Phi\!\left(z_\alpha - \dfrac{\Delta'-\Delta_0}{\sigma_{\bar{X}-\bar{Y}}}\right)$\\[3ex]
$H_a:\mu_1-\mu_2<\Delta_0$   & $1-\Phi\!\left(-z_\alpha - \dfrac{\Delta'-\Delta_0}{\sigma_{\bar{X}-\bar{Y}}}\right)$\\[3ex]
$H_a:\mu_1-\mu_2\ne\Delta_0$ & $\Phi\!\left(z_{\alpha/2} - \dfrac{\Delta'-\Delta_0}{\sigma_{\bar{X}-\bar{Y}}}\right) - \Phi\!\left(-z_{\alpha/2} - \dfrac{\Delta'-\Delta_0}{\sigma_{\bar{X}-\bar{Y}}}\right)$\\
\bottomrule
\end{tabular}
\end{center}

\textbf{Exercise.} Find the sample size for the one-tailed and two-tailed procedures.
\begin{itemize}
  \item For the one-tailed sample size, use $\beta(\Delta')=\beta$.
  \item For the two-tailed sample size formula, replace $\alpha$ with $\alpha/2$ in the one-tailed formula.
\end{itemize}
\emph{Answer} (with $m=n$): $\displaystyle n=\frac{(\sigma_1^2+\sigma_2^2)(z_\alpha+z_\beta)^2}{(\Delta'-\Delta_0)^2}$ for a one-tailed test, and the same with $z_{\alpha/2}$ in place of $z_\alpha$ for a two-tailed test.

\subsection*{10.2 The Two-Sample $t$ CI and Test}
\addcontentsline{toc}{subsection}{10.2 The Two-Sample $t$ CI and Test}

\textbf{Assumptions.} (1) The population distributions are normal; (2) $\sigma_1$ and $\sigma_2$ are unknown.

Previously we introduced the random variable
\[
Z = \frac{\bar{X}-\bar{Y}-(\mu_1-\mu_2)}{\sqrt{\dfrac{\sigma_1^2}{m}+\dfrac{\sigma_2^2}{n}}}.
\]
To get rid of the unknown $\sigma_1$ and $\sigma_2$ above, we can replace them with their sample standard deviations $S_1$ and $S_2$.

\textbf{Welch's theorem.} For two normal populations, the random variable
\[
T = \frac{\bar{X}-\bar{Y}-(\mu_1-\mu_2)}{\sqrt{\dfrac{S_1^2}{m}+\dfrac{S_2^2}{n}}}
\]
has approximately a $t$ distribution with df $\nu$ estimated by
\[
\nu = \frac{\left(\dfrac{s_1^2}{m}+\dfrac{s_2^2}{n}\right)^2}{\dfrac{(s_1^2/m)^2}{m-1}+\dfrac{(s_2^2/n)^2}{n-1}}
    = \frac{\left[(se_1)^2+(se_2)^2\right]^2}{\dfrac{(se_1)^4}{m-1}+\dfrac{(se_2)^4}{n-1}},
\]
where $se_1=s_1/\sqrt{m}$ and $se_2=s_2/\sqrt{n}$. This is called \textbf{Welch's df} (round $\nu$ down to an integer).

Then
\[
P\!\left(-t_{\alpha/2,\nu} < T = \frac{\bar{X}-\bar{Y}-(\mu_1-\mu_2)}{\sqrt{S_1^2/m+S_2^2/n}} < t_{\alpha/2,\nu}\right) \approx 1-\alpha .
\]

\textbf{The two-sample CI.} The two-sample CI for $\mu_1-\mu_2$ with confidence level approximately $100(1-\alpha)\%$ is
\[
\bar{x}-\bar{y} \pm t_{\alpha/2,\nu}\sqrt{\frac{s_1^2}{m}+\frac{s_2^2}{n}},
\]
where $\nu$ is Welch's df.

\textbf{Test procedure (two-sample $t$ test).} Null hypothesis $H_0:\mu_1-\mu_2=\Delta_0$. Test statistic value:
\[
t = \frac{\bar{x}-\bar{y}-\Delta_0}{\sqrt{\dfrac{s_1^2}{m}+\dfrac{s_2^2}{n}}}.
\]

\begin{center}
\begin{tabular}{ll}
\toprule
Alternative hypothesis & Rejection region\\
\midrule
$H_a:\mu_1-\mu_2>\Delta_0$   & $t\ge t_{\alpha,\nu}$ (upper-tailed)\\
$H_a:\mu_1-\mu_2<\Delta_0$   & $t\le -t_{\alpha,\nu}$ (lower-tailed)\\
$H_a:\mu_1-\mu_2\ne\Delta_0$ & either $t\ge t_{\alpha/2,\nu}$ or $t\le -t_{\alpha/2,\nu}$ (two-tailed)\\
\bottomrule
\end{tabular}
\end{center}
The $P$-value can be written similarly to the one-sample $t$ test.

\textbf{Large-sample procedure.}
\begin{itemize}
  \item If $m$ and $n$ are large ($m>40$ and $n>40$), by the CLT we get an approximately normal distribution.
  \item If $m$ and $n$ are large, then $\nu\ge\min(m-1,n-1)$ is also large. Then we can replace the $t$ critical value with the $z$ critical value to get an approximation to the $t$ procedure.
\end{itemize}

\textbf{Pooled $t$ procedures.} Assume $\sigma_1=\sigma_2=\sigma$ (even if it is unknown). Then we have the random variable
\[
Z = \frac{\bar{X}-\bar{Y}-(\mu_1-\mu_2)}{\sqrt{\dfrac{\sigma^2}{m}+\dfrac{\sigma^2}{n}}} = \frac{\bar{X}-\bar{Y}-(\mu_1-\mu_2)}{\sigma\sqrt{\dfrac{1}{m}+\dfrac{1}{n}}}.
\]
To get rid of $\sigma$ we could:
\begin{enumerate}
  \item replace $\sigma$ with $S_1$;
  \item replace $\sigma$ with $S_2$;
  \item use the ordinary average $(S_1+S_2)/2$ (very naive);
  \item best choice: use the \textbf{pooled} (i.e.\ combined) estimator of $\sigma^2$,
  \[
  S_p^2 = \frac{m-1}{m+n-2}S_1^2 + \frac{n-1}{m+n-2}S_2^2,
  \qquad\text{for which}\qquad
  \frac{(m+n-2)S_p^2}{\sigma^2}\sim\chi^2_{m+n-2}\ \ (\text{exercise}).
  \]
\end{enumerate}
The random variable is
\[
T = \frac{\bar{X}-\bar{Y}-(\mu_1-\mu_2)}{S_p\sqrt{\dfrac{1}{m}+\dfrac{1}{n}}} \sim t_{m+n-2}.
\]

\textbf{Pooled $t$ CI.} The pooled $t$ CI for $\mu_1-\mu_2$ with confidence level $100(1-\alpha)\%$ is
\[
\bar{x}-\bar{y}\pm t_{\alpha/2,m+n-2}\cdot s_p\sqrt{\frac{1}{m}+\frac{1}{n}} .
\]

\textbf{Pooled $t$ test.} Null hypothesis $H_0:\mu_1-\mu_2=\Delta_0$. Test statistic value:
\[
t = \frac{\bar{x}-\bar{y}-\Delta_0}{s_p\sqrt{\dfrac{1}{m}+\dfrac{1}{n}}},
\]
with rejection regions as before but with $m+n-2$ df.

\begin{itemize}
  \item When $\sigma_1=\sigma_2$, the pooled $t$ procedure is stronger than the two-sample (Welch's) $t$ procedure.
  \item When $\sigma_1\ne\sigma_2$, using the pooled $t$ gives incorrect results.
  \item When $\sigma_1\ne\sigma_2$ (or we are not sure), use the two-sample $t$ procedure unless there is strong evidence for using the pooled $t$ procedure.
\end{itemize}

\textbf{Power and type II errors, $\beta(\Delta')$.}
\begin{itemize}
  \item It is difficult to find the exact power and type II error probability.
  \item When $m$ and $n$ are large, we can use the formulas obtained for the two-sample $z$ procedure; in all those formulas, replace the $\sigma$'s with their estimates.
\end{itemize}

\subsection*{10.3 Analysis of Paired Data}
\addcontentsline{toc}{subsection}{10.3 Analysis of Paired Data}

\textbf{Assumptions.} Suppose the data consist of $n$ independently selected pairs
$(X_1,Y_1)$, $(X_2,Y_2)$, \dots, $(X_n,Y_n)$ with $E(X_i)=\mu_1$ and $E(Y_i)=\mu_2$.
Let $D_i=X_i-Y_i$. The $D_i$'s are assumed to be normally distributed with mean $\mu_D$ and standard deviation $\sigma_D$.
\begin{itemize}
  \item The $D_i$'s are independent of each other.
  \item $X$ and $Y$ within a pair are \emph{not} independent.
\end{itemize}
Consider $D=X-Y$. Then $E(D)=E(X-Y)=E(X)-E(Y)=\mu_1-\mu_2$, so $\mu_D=\mu_1-\mu_2$, and a CI for $\mu_D$ is equivalent to a CI for $\mu_1-\mu_2$.
Let $\bar{D}$ be the sample mean and $S_D$ the sample standard deviation of the differences. Then
\[
T = \frac{\bar{D}-\mu_D}{S_D/\sqrt{n}} \sim t_{n-1}.
\]

\textbf{Paired CI.} The paired $t$ CI for $\mu_D=\mu_1-\mu_2$ with confidence level $100(1-\alpha)\%$ has endpoints
\[
\bar{d}\pm t_{\alpha/2,n-1}\cdot\frac{s_D}{\sqrt{n}},
\]
where $\bar{d}$ and $s_D$ are the observed values. The one-sided confidence bounds are
$\bar{d}-t_{\alpha,n-1}\,s_D/\sqrt{n}$ (lower) and $\bar{d}+t_{\alpha,n-1}\,s_D/\sqrt{n}$ (upper).

\textbf{The paired $t$ test.} Null hypothesis $H_0:\mu_D=\Delta_0$. Test statistic value: $t=\dfrac{\bar{d}-\Delta_0}{s_D/\sqrt{n}}$.

\begin{center}
\begin{tabular}{ll}
\toprule
Alternative hypothesis & Rejection region\\
\midrule
$H_a:\mu_D>\Delta_0$   & $t\ge t_{\alpha,n-1}$ (upper-tailed)\\
$H_a:\mu_D<\Delta_0$   & $t\le -t_{\alpha,n-1}$ (lower-tailed)\\
$H_a:\mu_D\ne\Delta_0$ & either $t\le -t_{\alpha/2,n-1}$ or $t\ge t_{\alpha/2,n-1}$ (two-tailed)\\
\bottomrule
\end{tabular}
\end{center}
The $P$-value is computed as in the earlier $t$ test procedures.

\textbf{Comparison of the two statistics.}
\[
\text{Paired } t:\ \ t=\frac{\bar{d}-\Delta_0}{s_D/\sqrt{n}}
\qquad\qquad
\text{Two-sample } t:\ \ t=\frac{\bar{x}-\bar{y}-\Delta_0}{\sqrt{\dfrac{s_1^2}{m}+\dfrac{s_2^2}{n}}}
\]

\subsection*{10.4 Inferences About Two Population Proportions}
\addcontentsline{toc}{subsection}{10.4 Inferences About Two Population Proportions}

Let $X\sim\text{Bin}(m,p_1)$ and $Y\sim\text{Bin}(n,p_2)$ be independent. Consider the parameter $p_1-p_2$.
The natural estimator of $p_1-p_2$ is $\hat{p}_1-\hat{p}_2$, where $\hat p_1=X/m$, $\hat p_2=Y/n$.
For large samples $\hat{p}_1$ and $\hat{p}_2$ have approximately normal distributions, so $\hat{p}_1-\hat{p}_2$ also has an approximately normal distribution.
\[
E(\hat{p}_1-\hat{p}_2) = E(\hat{p}_1)-E(\hat{p}_2) = E\!\left(\frac{X}{m}\right)-E\!\left(\frac{Y}{n}\right) = \frac{E(X)}{m}-\frac{E(Y)}{n} = \frac{mp_1}{m}-\frac{np_2}{n} = p_1-p_2,
\]
\[
V(\hat{p}_1-\hat{p}_2) = V(\hat{p}_1)+V(\hat{p}_2) = \frac{V(X)}{m^2}+\frac{V(Y)}{n^2} = \frac{mp_1q_1}{m^2}+\frac{np_2q_2}{n^2} = \frac{p_1q_1}{m}+\frac{p_2q_2}{n},
\]
where $q_1=1-p_1$, $q_2=1-p_2$. So
$\sigma_{\hat{p}_1-\hat{p}_2}=\sqrt{\dfrac{p_1q_1}{m}+\dfrac{p_2q_2}{n}}$, and we have the random variable
\[
Z = \frac{\hat{p}_1-\hat{p}_2-(p_1-p_2)}{\sqrt{\dfrac{p_1q_1}{m}+\dfrac{p_2q_2}{n}}}. \tag{1}
\]

\textbf{A large-sample test procedure.} Let $H_0:p_1-p_2=0$ be the null hypothesis.
Let $p$ denote the common proportion, i.e.\ $p_1=p_2=p$, and $q=1-p$. Then the random variable (1) becomes
\[
Z = \frac{\hat{p}_1-\hat{p}_2-0}{\sqrt{pq\left(\dfrac{1}{m}+\dfrac{1}{n}\right)}}. \tag{2}
\]
To get rid of $p$ in the denominator, we can use the sample proportion as an estimate. The estimator of $p$ is
\[
\hat{p} = \frac{X+Y}{m+n} = \frac{X}{m+n}+\frac{Y}{m+n} = \frac{m}{m+n}\hat{p}_1+\frac{n}{m+n}\hat{p}_2, \tag{3}
\]
which is called the \emph{weighted average} of the estimators. Substituting (3) into (2), we obtain
\[
Z^* = \frac{\hat{p}_1-\hat{p}_2}{\sqrt{\hat{p}\hat{q}\left(\dfrac{1}{m}+\dfrac{1}{n}\right)}},\qquad \hat{q}=1-\hat{p}.
\]

\textbf{Two-proportion $z$ test.} Null hypothesis $H_0:p_1-p_2=0$. Test statistic value:
\[
z = \frac{\hat{p}_1-\hat{p}_2}{\sqrt{\hat{p}\hat{q}\left(\dfrac{1}{m}+\dfrac{1}{n}\right)}}.
\]

\begin{center}
\begin{tabular}{ll}
\toprule
Alternative hypothesis & Rejection region\\
\midrule
$H_a:p_1-p_2>0$   & $z\ge z_\alpha$\\
$H_a:p_1-p_2<0$   & $z\le -z_\alpha$\\
$H_a:p_1-p_2\ne0$ & either $z\le -z_{\alpha/2}$ or $z\ge z_{\alpha/2}$\\
\bottomrule
\end{tabular}
\end{center}
The $P$-value is computed as in the previous $z$ procedures.
These procedures are valid when $m\hat{p}_1\ge10$, $m\hat{q}_1\ge10$, $n\hat{p}_2\ge10$ and $n\hat{q}_2\ge10$.

\textbf{Power, $\beta$, and sample size determination.}
Because of the definition of $\beta$, we cannot use the random variables (2) and (3), as they assume $H_0$ is true; in the definition of $\beta$, $H_0$ is false.
$\beta$ is a function of both $p_1$ and $p_2$, i.e.\ $\beta(p_1,p_2)$ (textbook p.~605). For example, for $H_a:p_1-p_2>0$,
\[
\beta(p_1,p_2) = \Phi\!\left(\frac{z_\alpha\sqrt{\bar{p}\bar{q}\left(\frac{1}{m}+\frac{1}{n}\right)}-(p_1-p_2)}{\sigma}\right),
\quad
\bar p=\frac{mp_1+np_2}{m+n},\ \bar q = 1-\bar p,\ \sigma=\sqrt{\frac{p_1q_1}{m}+\frac{p_2q_2}{n}} .
\]
Setting $\beta(p_1,p_2)=\beta$ and solving gives the sample size.

\textbf{A large-sample two-proportion CI.} The CI for $p_1-p_2$, using the random variable (1), is
\[
\hat{p}_1-\hat{p}_2 \pm z_{\alpha/2}\sqrt{\frac{\hat{p}_1\hat{q}_1}{m}+\frac{\hat{p}_2\hat{q}_2}{n}} .
\]
This CI sometimes deviates from its nominal level. To overcome this, add one success and one failure to each sample;
then the confidence level is approximately correct. This adjusted CI can even be used for small samples.

\textbf{Small-sample procedures.}
\begin{itemize}
  \item Fisher's exact test.
  \item Add one success and one failure (adjusted CI).
  \item Use a software package to simulate the binomial distributions.
\end{itemize}

\subsection*{10.5 Inferences About Two Population Variances}
\addcontentsline{toc}{subsection}{10.5 Inferences About Two Population Variances}

\textbf{Theorem.} Let $X_1,\dots,X_m$ and $Y_1,\dots,Y_n$ be independent random samples from normal distributions with standard deviations $\sigma_1$ and $\sigma_2$.
Let $S_1$ and $S_2$ be the sample standard deviations. Then the random variable
\[
F = \frac{S_1^2/\sigma_1^2}{S_2^2/\sigma_2^2}
\]
has an $F$ distribution with $\nu_1=m-1$ and $\nu_2=n-1$ degrees of freedom.

\textbf{The $F$ test for equality of variances.} Null hypothesis $H_0:\sigma_1=\sigma_2$ ($\sigma_1^2=\sigma_2^2$). Test statistic value: $f=s_1^2/s_2^2$.

\begin{center}
\begin{tabular}{ll}
\toprule
Alternative hypothesis & Rejection region\\
\midrule
$H_a:\sigma_1>\sigma_2$   & $f\ge F_{\alpha,\nu_1,\nu_2}$\\
$H_a:\sigma_1<\sigma_2$   & $f\le F_{1-\alpha,\nu_1,\nu_2}$\\
$H_a:\sigma_1\ne\sigma_2$ & either $f\le F_{1-\alpha/2,\nu_1,\nu_2}$ or $f\ge F_{\alpha/2,\nu_1,\nu_2}$\\
\bottomrule
\end{tabular}
\end{center}
\textbf{$P$-value calculation:} use software.

\textbf{Exercise.} Derive the CI for $\sigma_1/\sigma_2$.

\emph{Answer.} From $P\big(F_{1-\alpha/2,\nu_1,\nu_2} < \frac{S_1^2/S_2^2}{\sigma_1^2/\sigma_2^2} < F_{\alpha/2,\nu_1,\nu_2}\big)=1-\alpha$, a $100(1-\alpha)\%$ CI for $\sigma_1^2/\sigma_2^2$ is
\[
\left(\frac{s_1^2}{s_2^2}\cdot\frac{1}{F_{\alpha/2,\nu_1,\nu_2}},\ \frac{s_1^2}{s_2^2}\cdot\frac{1}{F_{1-\alpha/2,\nu_1,\nu_2}}\right),
\]
and the CI for $\sigma_1/\sigma_2$ is obtained by taking square roots of the endpoints.
